\chapter*{Summary}
\addcontentsline{toc}{chapter}{Summary}
\markboth{Summary}{Summary}

\section*{What this document is}

DSTM \cite{dstm2026} is a distributed programming model for parallel
discrete-event and agent-based simulation. It provides a partitioned
distributed array versioned over simulation time, so producers and consumers
can reference the same state at different points on a shared timeline.
\Glspl{process} are thereby free to \gls{drift} out of lockstep without
rollback or global \glspl{barrier}.

The published description is high-level. This document supplies the ground
underneath it: what the objects are, what a program must guarantee, what the
runtime guarantees in return, and what follows from the combination.

It is \emph{prescriptive}: it specifies the model DSTM is meant to realise,
and records every departure of the current implementation in
\cref{ch:deviations}. Fourteen are recorded, six affecting soundness.

\section*{The central idea, in one sentence}

\begin{center}
\begin{minipage}{0.9\textwidth}
\itshape
DSTM is dynamic single assignment over the index space
$\Ranks \times \Time \times \Addr$, in which the time coordinate acts as a
renaming dimension that eliminates every anti- and output-dependence across
timesteps, and the \gls{comm-graph} supplies the per-channel \gls{lookahead} that makes
the residual flow-dependence graph acyclic --- and therefore schedulable
without rollback and without \glspl{barrier}.
\end{minipage}
\end{center}

Each mechanism then has a standard counterpart in dependence theory and
scheduling, and the properties follow from familiar results.
\Cref{tab:sync-dictionary} gives the dictionary; the essentials are that
versioning is renaming (\cref{thm:renaming}), a \gls{delta} is an edge weight
(\cref{lem:edge-span}), \gls{drift} is slack (\cref{def:slack}), and the
absence of a per-tick \gls{barrier} makes the cost a critical path rather than
a sum of per-tick maxima (\cref{thm:makespan}).

\section*{Structure}

\Cref{ch:definitions} fixes the vocabulary, separating the four quantities
called ``time'' --- origin, delta, effective, clock --- and the three called
``memory'' --- region, version, cell. The memory model
(\cref{sec:memory-model}) and execution model (\cref{sec:execution-model})
follow.

\Cref{ch:semantics} gives the semantics in three layers. $\Lzm$
(\cref{sec:l0}) is axiomatic: eight restrictions and one equation defining the
store. $\Lom$ (\cref{sec:l1}) is operational, its rules annotated with the
source lines implementing them. \Cref{sec:sync} recasts the system as a
dependence graph and derives the scheduling and memory results.

\Cref{ch:apis} presents the public interface against those objects, one
contract per entry point, separating checked preconditions from unchecked
program restrictions.

\section*{The restrictions}

DSTM offers an unusually strong memory model --- a read returns exactly the
byte written, with no coherence protocol, fences or visibility ordering. That
is paid for by eight restrictions on the program. Two carry most of the weight
and appear nowhere in the existing documentation or source:

\begin{center}
\begin{minipage}{0.88\textwidth}
\itshape
The \gls{lag} $\lag{p}{q}$ is a lower bound on how long a write takes to
become visible (\cref{res:lookahead}: $\delta \ge \lag{p}{q}$), and an upper
bound on how far back a read may look (\cref{res:staleness}:
$|\delta| \le \lag{p}{q}$).
\end{minipage}
\end{center}

This duality makes the paper's \glslink{pull-model}{pull} and
\glslink{push-model}{push} idioms both legal, and --- perhaps surprisingly ---
the same schedule (\cref{thm:pipeline-depth}). \Cref{ch:restrictions} lists
all eight with proposed checks; seven cost $O(1)$ per \gls{operation}.

\section*{Results}

\begin{table}[htbp]
\centering
\small
\begin{tabularx}{\textwidth}{@{}l >{\raggedright\arraybackslash}X l@{}}
\toprule
& \textbf{Result} & \textbf{Where} \\
\midrule
Safety & A published version is complete: no legal operation can still write
  to it & \cref{thm:write-closure} \\
       & Published contents never change & \cref{thm:immutability} \\
       & A retired version can never be legally read again
  & \cref{thm:discard-safety} \\
       & Every read returns exactly $\StoreH$ ($\Lom \refines \Lzm$)
  & \cref{thm:read-soundness} \\
\midrule
Determinacy & The result is independent of scheduling, message order and drift
  & \cref{thm:determinism} \\
            & \textbf{A DSTM execution computes exactly what the
  barrier-synchronised execution would} & \cref{thm:bsp} \\
\midrule
Liveness & Every read completes; deadlock freedom is equivalent to the absence
  of a zero-weight cycle in the \gls{comm-graph} & \cref{thm:deadlock-freedom} \\
         & The ring barrier establishes quiescence, and three rounds are
  necessary & \cref{thm:barrier-correct}, \cref{thm:barrier-three-rounds} \\
\midrule
Structure & Operations at distinct effective times are Bernstein-independent
  & \cref{thm:renaming} \\
          & Only flow dependences survive; the execution is in DSA form
  & \cref{cor:dsa} \\
          & The dependence graph is acyclic & \cref{thm:acyclic} \\
\midrule
Cost & Makespan is the critical path, versus a sum of per-tick maxima for BSP
  & \cref{thm:makespan} \\
     & Live versions are linear in drift and independent of run length
  & \cref{cor:space-drift} \\
     & Push and pull yield the same pipeline depth
  & \cref{thm:pipeline-depth} \\
\bottomrule
\end{tabularx}
\caption{The results proved in this document.}
\label{tab:results}
\end{table}

\Cref{thm:bsp} licenses the whole design. Removing the per-tick \gls{barrier}
and letting \glspl{process} \gls{drift} changes when work happens and how much
memory is live, provably not what is computed. Unlike optimistic simulation
nothing is ever discarded, because nothing is ever wrong.

\section*{Principal findings about the implementation}

Three findings are worth reading first.

\begin{enumerate}[leftmargin=1.5em,itemsep=4pt]

\item \textbf{The \gls{delta}/\gls{lag} contract is undocumented and unchecked}
  (\cref{dev:lookahead-uncheck}). \Cref{res:lookahead,res:staleness} are the
  assumptions that make both promotion guards sound, yet they appear in no
  header, comment or paper. Violating the first aborts; violating the second
  hangs, silently, inside a \gls{get}.

\item \textbf{A lag of zero is unsound, not merely unoptimised}
  (\cref{dev:zero-distance}). The paper states that ``a distance of 0 between
  all \glspl{rank} is a correct setup''. The implementation accepts such a matrix and
  then treats those pairs as \emph{unreachable}, suppressing all \glspl{heartbeat}
  and all visibility checks between them. Under the intended reading the
  configuration deadlocks (\cref{rem:zero-cycle}); as implemented it silently
  loses updates. The intended meaning --- all-to-all synchronisation every
  tick --- is expressed by a lag of one.

\item \textbf{\Glspl{version} do not inherit} (\cref{dev:non-inheriting}). Describing
  versioning as retaining ``a \gls{history} of state'' invites the reading that
  version $\eff$ is a snapshot of an evolving array. It is not: each version
  is zero-filled and independent (\cref{conv:no-inherit}). A program that
  writes at $\eff$ and reads the same address at $\eff+1$ without rewriting it
  reads zero. This is the most likely first-use error and it is documented
  nowhere.

\end{enumerate}

The remaining ten entries of \cref{ch:deviations} cover an unenforced
write-disjointness rule, which is the one violation that silently destroys
reproducibility (\cref{dev:write-disjoint}); asymmetric \glspl{distance-matrix} that
are only half-checked (\cref{dev:asymmetric}); two integer-width defects; a
missing bounds check; a non-collective error path; and three documentation
mismatches.
